% answers_short.tex —— 2019 年·上海春季招生数学试卷 · 简短答案速查
% 由 tools/make_short_answers.py 自 answers.tex 的 \gksolhead/\gkanswer 提取，
% 并与 2019/2019_数学真题_春季上海_答案要点.md 逐题比对后生成，勿手改（改源 MD 后重跑）。

\begin{gknotice}
\textbf{使用说明}：本节只给最终答案便于快速核对，完整推导见解析版。答案为编写者独立推导，\textbf{非官方标准答案}。
\end{gknotice}

\begin{gknotice}
选择题答案中的 U 表示待核，未给出猜测答案，详情见解析版。
\end{gknotice}

\gkpart{填空题}
\begin{itemize}[leftmargin=2.6em]
  \item[\textbf{1.}] $\{3,5\}$。
  \item[\textbf{2.}] $2$。
  \item[\textbf{3.}] $(-6,4)$。
  \item[\textbf{4.}] $y=\sqrt{x}$（$x>0$）。
  \item[\textbf{5.}] $2\sqrt2$。
  \item[\textbf{6.}] $a=-2$。
  \item[\textbf{7.}] $15$。
  \item[\textbf{8.}] $\sqrt{10}$。
  \item[\textbf{9.}] 若五天每天都须有人参加，为 $24$ 种；若允许同日多人及空日，为 $500$ 种。固定副本未显式说明每日要求，待核。
  \item[\textbf{10.}] $a=\sqrt3$。
  \item[\textbf{11.}] 若 $F_1,F_2$ 是椭圆两焦点，夹角范围为 $[\pi-\arccos(1/3),\pi]$；固定副本题面未明示两点定义，待核。
  \item[\textbf{12.}] $t=1$ 或 $t=-3$。
\end{itemize}

\gkpart{单选题}
\begin{center}\begin{tabular}{ccccc}
\toprule
题号 & 13 & 14 & 15 & 16 \\
\midrule
答案 & B & C & D & U \\
\bottomrule
\end{tabular}\end{center}

\gkpart{解答题}
\begin{itemize}[leftmargin=2.6em]
  \item[\textbf{17.}] (1) $\arccos(\sqrt3/4)$；(2) $\dfrac34$
  \item[\textbf{18.}] (1) $S_n=2n^2+n$；(2) 依各项非零的等比数列定义，$q\in(-1,0)\cup(0,\frac34)$；若采用允许零公比的递推约定，还须包含 $q=0$
  \item[\textbf{19.}] (1) $\dfrac83$；(2) $a=2$；(3) 若三点互异，则 $d(P_1)+d(P_3)>2d(P_2)$；若允许 $P_1=P_3$，大小关系不固定，题意待核
  \item[\textbf{20.}] (1) $\{-\dfrac{\sqrt3}{2},0,\dfrac{\sqrt3}{2}\}$；(2) $d=\dfrac{2\pi}{3}$ 或 $d=\pi$；(3) $T=3,4,5,6$
\end{itemize}
